\SourceLeaf{192}actual untying of the Gordian knot. It may be granted, however, that this omission
is of little consequence here, since the author's way of ruling out the
third case, if admissible, would carry over directly to this case as well.
Yet it is precisely this alleged proof of the impossibility of the third case that we cannot
accept as valid. The author presents the matter as follows. If, for example, in the third case
the first two segments are assumed to be equal and the third greater, then
$DC$ would be greater than $CL$. Now since $AML$ is likewise an isosceles triangle
to which the same basic condition applies as to the original triangle $AON$, then if
the third case with its assumed subdivision were the valid one, $DC=CL$ would have to hold,
contrary to the earlier finding.

We believe we have indicated at this point in the proof what matters in the author's view somewhat
more clearly and precisely than he himself did,
which also seems to us to make its weakness easier to recognize. For
it is plainly assumed quite arbitrarily that for all isosceles triangles with
angle $A$ at the vertex and a greater angle at the base, when the construction indicated in the theorem
is performed, the sequence of the segments cut off, as regards
their staying equal, increasing or decreasing, must in every case, regardless of the lengths
of the sides, necessarily be the same. Such an assumption surely cannot be regarded as self-evident.
But since the attempted proof of the impossibility
of the third case (and likewise of the fourth and fifth) rests on this alone, and the whole article
offers no other means of proving the impossibility of the case that was passed over,
we believe we have sufficiently justified the judgement expressed above,
while gladly giving full credit to the rest of the ingenious treatment in the following
articles.

\begin{center}\rule{3cm}{0.4pt}\end{center}

\Needspace{10\baselineskip}
\begin{center}\textsc{Remark.}\end{center}

{\small
This review, already printed in Volume IV, pages 368 to 370, has been reproduced here
for completeness. The figure, which appears in neither the original nor the first
reprinting, is intended to make the text easier to understand.

\begin{flushright}\textsc{Stäckel.}\end{flushright}
}

\begin{center}\rule{3cm}{0.4pt}\end{center}
